% ============================================================================
% 1. Introducción y recordatorio de self-attention
% ============================================================================
\section{Introducción}
\label{sec:intro}

Los modelos basados en \textit{Transformers} alcanzan un alto rendimiento en el procesamiento de secuencias, 
aunque su coste computacional aumenta considerablemente con la longitud de estas. Como alternativa, han 
surgido arquitecturas basadas en modelos de espacio de estados (\textit{State Space Models}, SSMs), que 
permiten procesar secuencias largas más eficientemente. Sin embargo, los SSMs tradicionales presentan 
limitaciones para seleccionar dinámicamente la información relevante en secuencias como las del lenguaje natural.

Para abordar esta limitación, Gu y Dao \cite{gu2023mamba} introdujeron Mamba, una arquitectura que incorpora 
un mecanismo de selección dependiente de la entrada. Este permite adaptar los parámetros en cada posición de 
la secuencia, controlando así qué información se incorpora al estado oculto y cuál se conserva del estado 
anterior.

En este trabajo, analizamos la dinámica de los tokens en Mamba a partir del estudio de Vo et al. \cite{vo2025token}. 
En particular, estudiamos los fenómenos de convergencia y divergencia de las representaciones, 
así como las condiciones y los factores que determinan su comportamiento.